\documentclass[a4paper,11pt]{ctexart}

\input{../mathnote-preamble.tex}

% 元信息配置
\renewcommand{\notetitle}{含摩擦的弹簧简谐运动问题}
\renewcommand{\notesubtitle}{模型建构·图像理解·解题策略}
\renewcommand{\noteauthor}{自学物理笔记}
\renewcommand{\notedate}{\today}
\renewcommand{\noteversion}{v1.0}

\MathNoteEnablePrint % 若需要送印前切换为 CMYK，可取消本行注释
\MathNoteEnableReviewStamp

\begin{document}

% 封面
\thispagestyle{empty}
\PageTag[secondary]{高中物理·简谐运动}
\setcounter{page}{1}

\noindent
\begin{minipage}[t][0.7\textheight][c]{\textwidth}
  \raggedleft
  \vspace*{\fill}
  {\sffamily\bfseries\Huge\color{accent}\notetitle\par}
  \vspace{0.6em}
  {\Large\color{inkgray}\notesubtitle\par}
  \vspace{0.5em}
  {\large\color{inkgray}\noteauthor\par}
  \vspace{0.4em}
  {\normalsize\color{inkgray}\noteversion\quad|\quad\notedate\par}
  \vspace*{\fill}
\end{minipage}

\begin{center}
\begin{minipage}{0.84\textwidth}
  \textcolor{inkgray}{\sffamily\small 学习导航}\\
  \begin{itemize}
    \item 从\keyword{物理情景}出发, 建立含摩擦的弹簧振动数学模型, 形成统一符号体系;
    \item 通过\inlinehint{受力分析}和\inlinehint{能量方法}, 推导振幅递减、停止条件等常考结论;
    \item 借助图像与表格, 总结可直接套用的\keyword{解题流程}与\keyword{题型模板}, 支持自学与复习。 
  \end{itemize}
\end{minipage}
\end{center}
\vspace{1em}

\clearpage
\thispagestyle{plain}
\phantomsection

% 目录
\tableofcontents
\newpage

% 正文开始

\section{情景引入与模型搭建}

\begin{definitionbox}{基本情景}
在光滑水平桌面上接一段弹簧, 一端固定于墙, 另一端连接一质量为 $m$ 的小物块。桌面并不完全光滑, 物块与桌面之间存在摩擦。将物块从平衡位置拉开一段距离 $A$, 松手后, 物块在弹簧和摩擦力作用下往复运动, 振幅逐渐减小, 最终停下。

本节研究的就是这类\keyword{含摩擦的弹簧简谐运动}问题, 重点放在:
\begin{itemize}
  \item 如何建立受力模型与数学方程;
  \item 振幅如何随时间或振动次数变化;
  \item 物块何时、在何处完全停止。
\end{itemize}
\end{definitionbox}

\subsection{物理模型与符号约定}

常用物理量及符号如表~\ref{tab:symbols} 所示。

\begin{table}[htbp]
  \centering
  \caption{常用物理量与符号约定}
  \label{tab:symbols}
  \begin{tabular}{ccl}
    \hline
    符号 & 单位 & 含义 \\
    \hline
    $m$ & $\mathrm{kg}$ & 物块质量 \\
    $k$ & $\mathrm{N/m}$ & 弹簧劲度系数 \\
    $\mu_s,\ \mu_k$ & 无 & 静、动摩擦因数 \\
    $x$ & $\mathrm{m}$ & 物块相对平衡位置的位移 (右正左负) \\
    $A$ & $\mathrm{m}$ & 振幅, 即位移的最大值 \\
    $T$ & $\mathrm{s}$ & 周期 \\
    $f$ & $\mathrm{Hz}$ & 频率, $f=1/T$ \\
    \hline
  \end{tabular}
\end{table}

为方便起见, 本笔记主要考虑\keyword{水平弹簧}模型; 竖直弹簧只需在受力分析中额外考虑重力与弹力平衡, 方法类似。

\subsection{含摩擦的水平弹簧示意图}

图~\ref{fig:model} 是一个带动摩擦的水平弹簧模型示意。习惯上取向右为位移、速度正方向。

\begin{figure}[htbp]
  \centering
  \begin{tikzpicture}[scale=1]
    % 地面和墙
    \draw[thick] (-0.5,0) -- (6,0);
    \fill[inkgray!40] (-0.5,0) rectangle (0,1.5);
    \draw[thick] (0,0) -- (0,1.5);
    % 弹簧
    \draw[decorate, decoration={coil, aspect=0.7, segment length=5pt, amplitude=3pt}] 
      (0,0.75) -- (3,0.75);
    % 物块
    \draw[fill=surface] (3,0.25) rectangle (4.3,1.25);
    \node at (3.65,0.75) {$m$};
    % 坐标与速度方向
    \draw[->, thick] (3.65,0.25) -- (5.1,0.25) node[right] {$x$};
    \draw[->, thick] (3.65,-0.1) -- (4.8,-0.1) node[right] {$v>0$};
    % 力
    \draw[->, thick, accent] (3.0,0.75) -- (2.0,0.75) node[left] {$F_\text{弹}$};
    \draw[->, thick, highlight] (3.65,0.25) -- (2.7,0.25) node[left] {$f$};
    \draw[->, thick] (3.65,1.25) -- (3.65,2.0) node[above] {$N$};
    \draw[->, thick] (3.65,0.25) -- (3.65,-0.6) node[below] {$mg$};
  \end{tikzpicture}
  \caption{含摩擦的水平弹簧振动模型示意}
  \label{fig:model}
\end{figure}

\begin{focuspoints}
  \item 摩擦力大小 $f$ 在滑动时为 $f=\mu_k N=\mu_k mg$, 方向总是\keyword{与相对运动方向相反};
  \item 静止时, 摩擦力 $f_s\le \mu_s N$, 自动调节, 用于维持\keyword{平衡或即将起动}状态;
  \item 含摩擦的弹簧振动问题, 本质是\keyword{在简谐运动基础上叠加摩擦做功, 使机械能逐渐减少}。
\end{focuspoints}

\section{无摩擦弹簧简谐运动快速回顾}

在研究含摩擦的复杂情形前, 先回顾理想弹簧的简谐运动结果, 后续许多结论都在此基础上做修正。

\begin{summarybox}{无摩擦水平弹簧振动要点}
考虑质量为 $m$ 的物块与劲度系数为 $k$ 的弹簧, 水平放置, 无摩擦:
\begin{itemize}
  \item 受力: 物块只受弹力 $F_\text{弹}=-kx$;
  \item 牛顿第二定律: $ma=-kx\Rightarrow \ddot{x}+\dfrac{k}{m}x=0$;
  \item 运动形式: 简谐运动
  \[
    x(t)=A\cos(\omega t+\varphi),\quad \omega=\sqrt{\dfrac{k}{m}};
  \]
  \item 周期与频率:
  \[
    T=\dfrac{2\pi}{\omega}=2\pi\sqrt{\dfrac{m}{k}},\qquad f=\dfrac{1}{T};
  \]
  \item 机械能守恒:
  \[
    E=\dfrac{1}{2}kx^2+\dfrac{1}{2}mv^2=\text{常量}.
  \]
\end{itemize}
\end{summarybox}

\begin{focuspoints}
  \item 理想情况下, 振幅 $A$ 与机械能 $E$ 保持不变;
  \item 引入摩擦后, 周期 $T$ 通常\keyword{变化不大}, 但振幅和机械能都会\keyword{逐渐减小};
  \item 因此, 处理含摩擦问题常常仍把单次往返看作“近似简谐运动”, 再通过能量方法研究振幅的变化。
\end{focuspoints}

\section{含滑动摩擦的弹簧振动}

本节考虑物块在运动过程中一直处于滑动状态, 摩擦力大小恒为 $f=\mu_k mg$, 方向总是与速度方向相反。

\subsection{受力分析与起振条件}

\begin{theorembox}{起振条件}
若将物块从平衡位置拉到一侧, 使弹簧伸长或压缩 $A$, 然后由静止释放, 则物块能够开始周期性往返运动的必要条件是
\[
  kA > \mu_s mg,
\]
即\keyword{最大弹力}大于\keyword{最大静摩擦力}。若 $kA\le \mu_s mg$, 物块释放后将\keyword{直接静止}在某个偏离平衡的位置。
\end{theorembox}

实际解题中, 常用下列步骤判断:

\begin{roadmap}
  \RoadmapStep{画出物块受力图, 包括弹力、摩擦力、重力和支持力;}
  \RoadmapStep{在最大偏离位置利用 $F_\text{弹,max}=kA$ 估算弹力上限;}
  \RoadmapStep{比较 $kA$ 与 $\mu_s mg$, 判断物块能否被“拉回”；}
  \RoadmapStep{若不能起振, 直接用静力平衡求出最终停留位置;}
  \RoadmapStep{若能起振, 则转入含动摩擦的振幅递减分析。}
\end{roadmap}

\subsection{振幅线性递减规律 (能量法)}

下面考虑物块已经开始周期性往返运动, 且在运动过程中始终处于滑动状态 (即静摩擦可以忽略)。这时摩擦力大小恒为 $f=\mu_k mg$。

\begin{theorembox}{振幅线性递减}
设物块在第 $n$ 次振动的最大位移(振幅)为 $A_n$, 则相邻两次振动的振幅满足
\[
  A_{n+1}=A_n-\Delta A,\qquad 
  \Delta A=\dfrac{2\mu_k mg}{k},
\]
即振幅以\keyword{等差数列}形式逐次减小, 每次减小量 $\Delta A$ 为常数。
\end{theorembox}

\begin{proofbox}{推导思路}
\begin{enumerate}
  \item 取相邻两次振动的右端极值点为研究对象, 设第 $n$ 次振幅为 $A_n$, 第 $n+1$ 次为 $A_{n+1}$。
  \item 在任一极值点速度为零, 故机械能全为弹性势能:
  \[
    E_n=\frac{1}{2}kA_n^2,\qquad E_{n+1}=\frac{1}{2}kA_{n+1}^2.
  \]
  \item 从 $+A_n$ 运动到 $-A_{n+1}$ 的过程中, 物块所经过的路程约为 $A_n+A_{n+1}$, 动摩擦力大小恒为 $f=\mu_k mg$, 方向与速度相反, 故摩擦力做功
  \[
    W_f = -f(A_n+A_{n+1})=-\mu_k mg(A_n+A_{n+1}).
  \]
  \item 由机械能变化
  \[
    E_{n+1}=E_n+W_f,
  \]
  得
  \[
    \frac{1}{2}kA_{n+1}^2=\frac{1}{2}kA_n^2-\mu_k mg(A_n+A_{n+1}).
  \]
  \item 两边乘以 $\dfrac{2}{k}$, 并整理得
  \[
    (A_n-A_{n+1})(A_n+A_{n+1})=\frac{2\mu_k mg}{k}(A_n+A_{n+1}).
  \]
  \item 由于 $A_n,A_{n+1}>0$, 可约去 $(A_n+A_{n+1})$, 得
  \[
    A_n-A_{n+1}=\dfrac{2\mu_k mg}{k}=\Delta A.
  \]
\end{enumerate}
因此 $A_{n+1}=A_n-\Delta A$, 证毕。
\end{proofbox}

由此可得振幅的一般表达式
\[
  A_n = A_0 - n\Delta A,\qquad n=0,1,2,\dots,
\]
其中 $A_0$ 为初始振幅。

\begin{table}[htbp]
  \centering
  \caption{振幅随振动次数的变化 (示意)}
  \label{tab:amplitude}
  \begin{tabular}{ccc}
    \hline
    振动序号 $n$ & 振幅 $A_n$ & 说明 \\
    \hline
    $0$ & $A_0$ & 初始拉开的振幅 \\
    $1$ & $A_0-\Delta A$ & 第一次往返后的振幅 \\
    $2$ & $A_0-2\Delta A$ & 第二次往返后的振幅 \\
    $3$ & $A_0-3\Delta A$ & 以此类推 \\
    $\vdots$ & $\vdots$ & $\vdots$ \\
    \hline
  \end{tabular}
\end{table}

\begin{notebox}{停止条件}
当振幅减小到一定程度时, 弹力的最大值 $kA_n$ 将不足以克服静摩擦力 $\mu_s mg$, 物块将停止振动并静止在某个偏离平衡的位置。停止条件可写成
\[
  kA_n \le \mu_s mg,
\]
于是最大的振动次数 $n_{\max}$ 约为满足
\[
  A_{n_{\max}} > \frac{\mu_s mg}{k},\quad 
  A_{n_{\max}+1} \le \frac{\mu_s mg}{k}
\]
的整数 $n_{\max}$。
\end{notebox}

\subsection{位移—时间图像 (示意)}

含摩擦时, 真实运动并非严格的简谐运动, 但若摩擦力相对较小, 可以把每一周期看作近似简谐运动, 且周期变化不大, 振幅随时间缓慢减小。图~\ref{fig:xt} 给出了一个示意图。

\begin{figure}[htbp]
  \centering
  \begin{tikzpicture}[scale=0.9]
    \draw[->] (0,0) -- (8,0) node[right] {$t$};
    \draw[->] (0,-2) -- (0,2) node[above] {$x$};
    % 上下包络线
    \draw[accent, thick] (0,1.6) -- (7,0.6);
    \draw[accent, thick] (0,-1.6) -- (7,-0.6);
    % 近似振动曲线 (示意用指数衰减包络)
    \draw[secondary, smooth, domain=0:7, samples=200]
      plot (\x, {1.6*exp(-0.15*\x)*cos(2*deg(\x))});
    \node[anchor=west] at (4.2,1.2) {\small 振幅缓慢减小};
  \end{tikzpicture}
  \caption{含摩擦弹簧振动的位移—时间示意图}
  \label{fig:xt}
\end{figure}

\begin{focuspoints}
  \item 理想情况下, 振幅为水平直线; 含摩擦时, 振幅随时间呈\keyword{递减趋势};
  \item 图像法常用于比较某时刻 $x,v,a$ 的大小与符号, 例如通过斜率判断速度方向;
  \item 在选择题中, 通过\keyword{振幅逐渐减小}和\keyword{能量不断损失}这两个特征, 往往可以快速排除错误选项。
\end{focuspoints}

\section{常见问题与解题思路整理}

含摩擦的弹簧振动题目在高中阶段常以\keyword{综合题}形式出现, 既考察力学基础, 也考察能量方法和图像理解。下面用一张表格总结常见问法与对应思路。

\begin{table}[htbp]
  \centering
  \caption{含摩擦弹簧问题的常见问法与解题思路}
  \label{tab:qa}
  \begin{tabular}{p{0.32\textwidth} p{0.55\textwidth}}
    \hline
    题目问法 & 解题思路 \\
    \hline
    能否开始振动 & 比较 $kA$ 与 $\mu_s mg$, 结合受力图判断能否克服最大静摩擦力。 \\
    第 $n$ 次振动末的振幅 & 利用 $A_n=A_0-n\Delta A$, 先求 $\Delta A=\dfrac{2\mu_k mg}{k}$, 再代入。 \\
    停止前经历的最大振动次数 $n_{\max}$ & 由 $kA_n>\mu_s mg,\ kA_{n+1}\le \mu_s mg$ 求出 $n_{\max}$, 注意取整数。 \\
    停止前的总路程 & 将每一往返路程 $4A_n$ 求和, 通常是一个等差数列求和问题。 \\
    比较不同摩擦条件下的“快慢” & 周期近似相同, 主要比较\keyword{停止时间}、\keyword{总路程}或\keyword{能量损失速率}。 \\
    \hline
  \end{tabular}
\end{table}

\begin{summarybox}{解题通用流程}
\begin{enumerate}
  \item \textbf{画图}: 画出物理情景与受力图, 标注位移、速度方向与各力;
  \item \textbf{起振判断}: 用 $kA$ 与 $\mu_s mg$ 判断能否起振, 若不能起振, 直接静力平衡求解;
  \item \textbf{能量分析}: 若能起振, 用机械能变化关系推导振幅递减规律 $A_{n+1}=A_n-\Delta A$;
  \item \textbf{数列处理}: 将振幅视为等差数列, 求 $A_n$, $n_{\max}$ 以及总路程等量;
  \item \textbf{图像检查}: 用位移—时间图或能量—时间图检查结果是否合理, 防止漏掉摩擦做功等因素。
\end{enumerate}
\end{summarybox}

\section{竖直弹簧模型拓展}

在许多竞赛题或高考压轴题中, 会把\keyword{竖直弹簧}与\keyword{重力}、\keyword{摩擦}结合考察。本节只做最核心、最常用的整理, 重点是用\keyword{平衡位置}统一处理竖直弹簧的振动问题。

\subsection{平衡位置与等效模型}

考虑劲度系数为 $k$ 的轻弹簧竖直悬挂, 下端挂质量为 $m$ 的小物块, 如图~\ref{fig:vertical-spring}。设弹簧原长为 $L_0$, 物块静止时弹簧总长为 $L$, 则平衡时有
\[
  k(L-L_0)=mg \quad\Rightarrow\quad L-L_0=\frac{mg}{k}.
\]
记平衡位置时弹簧被拉长的长度为
\[
  x_0 = \frac{mg}{k}.
\]

\begin{figure}[htbp]
  \centering
  \begin{tikzpicture}[scale=1]
    % 悬点
    \fill (0,3) circle (1pt);
    \draw[thick] (-0.5,3.2) -- (0.5,3.2);
    \draw[thick] (0,3.2) -- (0,3);
    % 弹簧
    \draw[decorate, decoration={coil, aspect=0.7, segment length=4pt, amplitude=3pt}]
      (0,3) -- (0,1.4);
    % 物块
    \draw[fill=surface] (-0.4,1.4) rectangle (0.4,0.6);
    \node at (0,1.0) {$m$};
    % 位移坐标
    \draw[->, thick] (0.8,1.0) -- (0.8,-0.2) node[right] {$y$};
    \node[anchor=west] at (0.8,1.0) {\small $y=0$ (平衡)};
    % 力
    \draw[->, thick, accent] (0,0.6) -- (0,-0.2) node[below] {$mg$};
    \draw[->, thick, secondary] (0,1.4) -- (0,2.2) node[above] {$F_\text{弹}$};
  \end{tikzpicture}
  \caption{竖直弹簧与平衡位置示意}
  \label{fig:vertical-spring}
\end{figure}

若令 $y$ 为物块相对\keyword{平衡位置}的位移 (下为正), 则
\[
  x = x_0 + y
\]
为弹簧相对\keyword{自然长度}的伸长量。此时弹力大小为 $F_\text{弹}=k(x_0+y)$, 方向向上, 重力为 $mg$, 方向向下。对物块列方程:
\[
  ma = \sum F = mg - k(x_0+y).
\]
利用平衡条件 $k x_0 = mg$, 上式化简为
\[
  ma = mg - kx_0 - ky = -ky,
\]
即
\[
  \ddot{y}+\frac{k}{m}y=0.
\]

\begin{summarybox}{竖直弹簧等效为水平弹簧}
\begin{itemize}
  \item 若以\keyword{平衡位置}为坐标原点 ($y=0$), 则竖直弹簧的运动方程与水平弹簧\keyword{完全相同}:
  \[
    \ddot{y}+\frac{k}{m}y=0;
  \]
  \item 因此, 周期、频率等结论与水平弹簧一致:
  \[
    T=2\pi\sqrt{\frac{m}{k}},\qquad f=\frac{1}{T};
  \]
  \item 竖直弹簧所有“新”现象, 本质上都是\keyword{平衡位置相对于自然长度发生了平移}, 可通过 $x_0=\dfrac{mg}{k}$ 统一处理。
\end{itemize}
\end{summarybox}

\subsection{竖直弹簧与摩擦的组合思路}

若物块沿竖直导轨上下滑动, 与导轨之间存在动摩擦力 $f=\mu_k N$, 则
\begin{itemize}
  \item 支持力 $N$ 通常与导轨的几何关系有关 (例如竖直导轨时 $N$ 为水平向, 与重力无关);
  \item 以平衡位置为原点, 仍可把弹力与重力的合力写成 $-ky$, 摩擦力则起到\keyword{额外消耗机械能}的作用, 处理方式与水平面上的摩擦类似;
  \item 振幅递减公式 $A_{n+1}=A_n-\Delta A$ 仍成立, 只需把 $\mu_k mg$ 替换为\keyword{实际摩擦力大小} $f=\mu_k N$。
\end{itemize}

\begin{notebox}{典型考点}
\begin{itemize}
  \item 先用静力平衡求出竖直弹簧的平衡伸长 $x_0$;
  \item 再以平衡位置为原点分析振动, 直接套用水平弹簧简谐运动与含摩擦振幅递减的结论;
  \item 注意区分“相对平衡位置的振幅”与“相对自然长度的最大伸长(压缩量)”。
\end{itemize}
\end{notebox}

\section{典型例题与解题模板}

\subsection{例题: 振幅递减与总路程}

\begin{examplebox}{含摩擦水平弹簧振动}
如图~\ref{fig:model} 所示, 质量为 $m$ 的小物块与劲度系数为 $k$ 的轻弹簧连接, 放在水平粗糙桌面上。已知物块与桌面的静、动摩擦因数分别为 $\mu_s,\mu_k$, 且 $\mu_s \approx \mu_k$。将物块从平衡位置向右缓慢拉至位移 $A_0$, 然后由静止释放。

\begin{enumerate}
  \item 判断物块能否开始周期性往返运动;
  \item 若可以, 写出第 $n$ 次振动末的振幅 $A_n$ 表达式;
  \item 求物块在完全停止前经过的总路程 $S$ 的表达式。
\end{enumerate}
\end{examplebox}

\begin{proofbox}{解析与解题模板}
\textbf{(1) 起振条件}

最大弹力 $F_{\text{弹,max}}=kA_0$, 最大静摩擦力 $F_{s,\max}=\mu_s mg$。若
\[
  kA_0 \le \mu_s mg,
\]
则物块释放后弹力始终不能战胜静摩擦力, 物块不会往返振动, 而是直接静止在某个偏移平衡的位置。若
\[
  kA_0 > \mu_s mg,
\]
则物块可以被拉回, 开始往返运动。

\medskip
\textbf{(2) 振幅递减公式}

在物块已经往返运动的阶段, 近似认为它始终处于滑动状态, 摩擦力大小恒为 $f=\mu_k mg$。由上一节结论直接得到
\[
  A_{n+1}=A_n-\Delta A,\qquad 
  \Delta A=\dfrac{2\mu_k mg}{k},
\]
于是
\[
  A_n = A_0 - n\Delta A
      = A_0 - n\,\dfrac{2\mu_k mg}{k}.
\]

\medskip
\textbf{(3) 停止前的总路程}

当某次振动后的振幅减小到
\[
  A_n \le \frac{\mu_s mg}{k}
\]
时, 最大弹力 $kA_n$ 已不足以克服静摩擦力, 物块将不再完成新的完整往返, 逐渐停在某个偏移位置。令
\[
  A_n = A_0 - n\Delta A > \frac{\mu_s mg}{k},\quad
  A_{n+1} = A_0 - (n+1)\Delta A \le \frac{\mu_s mg}{k},
\]
即可求出最大的完整往返次数 $n_{\max}$。

在前 $n_{\max}$ 次完整往返中, 第 $n$ 次往返的路程约为 $4A_n$ (从右极值到左极值再回到右极值), 因此总路程为
\[
  S \approx \sum_{n=0}^{n_{\max}} 4A_n
   = 4\sum_{n=0}^{n_{\max}}\bigl(A_0 - n\Delta A\bigr),
\]
这是一个等差数列求和问题:
\[
  S = 4\left[(n_{\max}+1)A_0 - \frac{n_{\max}(n_{\max}+1)}{2}\Delta A\right].
\]

\begin{focuspoints}
  \item 整题结构为“起振判断 $\rightarrow$ 振幅数列 $\rightarrow$ 等差求和”;
  \item 只要记住 $\Delta A=\dfrac{2\mu_k mg}{k}$, 大部分相关题目都可按同一模板处理;
  \item 若题目给出具体数值, 只需在最后一步代入计算即可。
\end{focuspoints}
\end{proofbox}

\subsection{例题: 数值综合——振幅、次数与路程}

\begin{examplebox}{含摩擦水平弹簧的数值计算}
质量为 $m=0.50\,\mathrm{kg}$ 的小物块与劲度系数为 $k=50\,\mathrm{N/m}$ 的轻弹簧连接, 放在水平粗糙桌面上。已知物块与桌面之间的静、动摩擦因数分别为 $\mu_s=0.20,\ \mu_k=0.15$, 取 $g=10\,\mathrm{m/s^2}$。将物块从平衡位置向右缓慢拉至 $A_0=0.30\,\mathrm{m}$ 处, 然后由静止释放。

\begin{enumerate}
  \item 判断物块能否开始周期性往返运动;
  \item 求每次完整往返后振幅的减少量 $\Delta A$;
  \item 求物块在完全停止前大约能完成多少次\keyword{完整往返} (从右端到左端再回到右端);
  \item 估算物块在完全停止前经过的总路程 $S$。
\end{enumerate}
\end{examplebox}

\begin{proofbox}{解析}
\textbf{(1) 起振判断}

最大弹力
\[
  F_{\text{弹,max}} = kA_0 = 50\times 0.30 = 15\,\mathrm{N},
\]
最大静摩擦力
\[
  F_{s,\max} = \mu_s mg = 0.20\times 0.50\times 10 = 1.0\,\mathrm{N}.
\]
因为 $F_{\text{弹,max}}\gg F_{s,\max}$, 故物块必然能够被拉回平衡位置并开始往返运动。

\medskip
\textbf{(2) 振幅减少量 $\Delta A$}

动摩擦力大小为
\[
  f = \mu_k mg = 0.15\times 0.50\times 10 = 0.75\,\mathrm{N}.
\]
由振幅递减公式
\[
  \Delta A = \frac{2f}{k} = \frac{2\times 0.75}{50} = 0.03\,\mathrm{m}.
\]

\medskip
\textbf{(3) 完整往返次数}

第 $n$ 次完整往返后的振幅
\[
  A_n = A_0 - n\Delta A = 0.30 - 0.03n.
\]
物块停止振动的条件是最大弹力不足以克服静摩擦力:
\[
  kA_n \le \mu_s mg 
  \quad\Rightarrow\quad 
  A_n \le \frac{\mu_s mg}{k}
  = \frac{0.20\times 0.50\times 10}{50}
  = 0.02\,\mathrm{m}.
\]
代入 $A_n=0.30-0.03n$ 得
\[
  0.30-0.03n \le 0.02
  \quad\Rightarrow\quad
  0.28 \le 0.03n
  \quad\Rightarrow\quad
  n \ge \frac{0.28}{0.03}\approx 9.3.
\]
因此, 当 $n=9$ 时仍有 $A_9\approx 0.30-0.27=0.03\,\mathrm{m}>0.02\,\mathrm{m}$, 还能完成第 $9$ 次完整往返; 当 $n=10$ 时 $A_{10}$ 已接近或小于静摩擦对应的最小振幅, 物块不再完成新的完整往返。故可认为大约完成
\[
  n_{\max}\approx 9
\]
次完整往返。

\medskip
\textbf{(4) 完全停止前的总路程}

在前 $n_{\max}$ 次完整往返中, 第 $n$ 次往返的路程约为 $4A_n$, 因此
\[
  S \approx \sum_{n=0}^{n_{\max}}4A_n
    = 4\sum_{n=0}^{9}(0.30-0.03n).
\]
利用等差数列求和,
\[
  \sum_{n=0}^{9}(0.30-0.03n)
  = 10\times 0.30 - 0.03\frac{9\times 10}{2}
  = 3.0 - 1.35
  = 1.65\,\mathrm{m},
\]
于是
\[
  S \approx 4\times 1.65 = 6.6\,\mathrm{m}.
\]

\begin{focuspoints}
  \item 数值题的关键在于先写出\keyword{通用公式}, 再代入具体数值, 避免在推导过程中就带入数字;
  \item 对 $n_{\max}$ 的判断一定要结合\keyword{不等式与取整}, 注意“还能再完成一次”和“刚好不能完成”的区别;
  \item 总路程的计算几乎总是转化为\keyword{等差数列求和}问题, 建议熟练掌握 $\displaystyle \sum_{n=0}^{N}(a+nd)$ 的计算。
\end{focuspoints}
\end{proofbox}

\section{自测练习与拓展思考}

\begin{notebox}{自测练习}
\begin{enumerate}
  \item 质量为 $0.50\,\mathrm{kg}$ 的小物块与劲度系数为 $50\,\mathrm{N/m}$ 的轻弹簧连接, 放在水平粗糙桌面上, 动摩擦因数 $\mu_k=0.10$, 静摩擦因数略大于 $0.10$。将物块从平衡位置拉到右侧 $0.20\,\mathrm{m}$ 处, 由静止释放。
  \begin{enumerate}
    \item 判断物块能否开始往返运动;
    \item 求每次完整往返后振幅减少量 $\Delta A$;
    \item 利用 $A_n=A_0-n\Delta A$ 估算物块完成多少次完整往返后基本停止。
  \end{enumerate}
  \item 某题中给出含摩擦弹簧振动的位移—时间图像, 发现相邻两个波峰的间隔几乎不变, 但波峰高度逐渐下降。试从能量角度解释这一现象, 并指出摩擦力在其中起到的作用。
  \item 尝试把本笔记中的水平弹簧模型改为竖直弹簧, 分析物块在竖直方向运动时的平衡位置、起振条件与振幅递减规律是否发生变化, 说明理由。
\end{enumerate}
\end{notebox}

\begin{summarybox}{本节小结}
\begin{itemize}
  \item 含摩擦的弹簧振动可视为“简谐运动 + 摩擦做功导致能量损失”的组合模型;
  \item 动摩擦力恒定时, 通过对相邻极值点应用机械能变化关系, 可得到\keyword{振幅线性递减}规律 $A_{n+1}=A_n-\dfrac{2\mu_k mg}{k}$;
  \item 利用振幅的等差数列特性, 可以系统解决“第 $n$ 次振幅”、停止前的最大振动次数、总路程等一类问题;
  \item 图像法和能量观点是检验答案合理性的重要工具, 建议在练习中主动使用。
\end{itemize}
\end{summarybox}

\end{document}
